% Konklusjon

\begin{frame}{Oppsummering}
\small
\begin{itemize}\setlength{\itemsep}{5pt}
  \item \textbf{Ingen generell jobbkrise.} Sysselsettingen i privat sektor vokser; sysselsettingsraten er stabil.
  \item \textbf{Hele fordelingen.} Siden oktober 2022 har sysselsettingen falt 0,1\,\% i de mest eksponerte yrkene mot 0,6\,\% i de minst eksponerte: et gap på $+0{,}5$ prosentpoeng, godt innenfor bootstrap-standardfeilen på rundt 2,5.
  \item \textbf{Kanarifuglen er reell, men generaliserer ikke.} Unge i programvareutvikling og kundeservice falt etter ChatGPT; på tvers av hele eksponeringsfordelingen falt ikke Q5 bak Q1 i 21--30. Pretrendene maner til varsomhet begge veier.
  \item \textbf{Landene imellom.} Norge ligner USA og Sverige i casestudiene, Danmark og Finland i den systematiske analysen.
  \item \textbf{Overvåkingen fortsetter} månedlig på kiindeksen.no, og utvidelser vurderes. Alle resultater gjelder privat sektor.
\end{itemize}
\end{frame}
